%%
%% XXII_universality_theorem.tex
%%
%% The Perturbative Universality Theorem for Paper XXII.
%%
%% This is the actual main conceptual result of the programme:
%% arithmetic structure survives only perturbatively in the Airy limit.
%%
%% NOT an organisational document. A mathematical statement.
%%
\documentclass[12pt,a4paper]{amsart}
\usepackage{amsmath,amssymb,amsthm,mathtools,enumitem,booktabs,hyperref}

\newtheorem{theorem}{Theorem}[section]
\newtheorem{lemma}[theorem]{Lemma}
\newtheorem{proposition}[theorem]{Proposition}
\newtheorem{corollary}[theorem]{Corollary}
\theoremstyle{definition}
\newtheorem{definition}[theorem]{Definition}
\newtheorem{remark}[theorem]{Remark}

\newcommand{\Ai}{\mathrm{Ai}}
\newcommand{\KAi}{K^{\mathrm{Airy}}}
\newcommand{\DAi}{D^{\mathrm{Airy}}}
\newcommand{\PAi}{\Pi_{\mathrm{Airy}}}
\newcommand{\Aarith}{A^{\mathrm{arith}}}
\newcommand{\AAiry}{A^{\mathrm{Airy}}}
\newcommand{\norm}[1]{\|#1\|}
\newcommand{\ip}[2]{\langle #1,\, #2\rangle}
\newcommand{\Ec}{E_c}

\title{The Perturbative Universality Theorem\\
{\normalsize Paper XXII: Arithmetic Sampling at the Airy Edge}}
\author{Sebastian Schmalnauer}
\date{May 2026}

\begin{document}
\maketitle

\begin{abstract}
We state and prove the Perturbative Universality Theorem for the
arithmetic sampling operator at the prolate spheroidal Airy edge.
The theorem asserts a precise scale separation:
the universal Airy gap $\delta_{\mathrm{Airy}} \sim O(1)$ dominates
the arithmetic perturbation $\|E_c\|_{\mathrm{op}} = o(1)$,
so that the spectral gap of the mismatch operator $\DAi$ persists
unconditionally for large $c$.
The arithmetic structure of the prime sampling sequence
survives the Airy scaling limit, but only perturbatively:
it cannot destroy the universal gap, though it modulates its value
at subleading order.
This is the precise formulation of universality for the XXII programme.
\end{abstract}

\tableofcontents

%%=================================================================
\section{The Scale Separation}
\label{sec:scales}
%%=================================================================

\subsection{Two scales in the Airy limit}

Let $c \to \infty$ and let $\DAi = \PAi(\Aarith - \KAi)\PAi$
be the Airy mismatch operator (Paper XXII, Section 5).
Decompose:
\begin{equation}\label{eq:decomp}
  \Aarith = \underbrace{\AAiry}_{= I} + \Ec,
  \qquad \Ec := \Aarith - I,
\end{equation}
so that:
\begin{equation}\label{eq:mismatch_decomp}
  \DAi
  = \underbrace{\PAi(I - \KAi)\PAi}_{\text{universal part}}
  + \underbrace{\PAi \Ec \PAi}_{\text{arithmetic part}}.
\end{equation}

\begin{definition}[Universal gap and arithmetic perturbation]
\begin{align}
  \delta_{\mathrm{Airy}}
  &:= \inf\sigma(\PAi(I - \KAi)\PAi)
   = 1 - \sup\sigma(\KAi|_{\PAi L^2})
   \ge F_{\mathrm{TW}}(0)
   \approx 0.9597,
  \label{eq:delta_def}\\
  \varepsilon_c
  &:= \norm{\PAi \Ec \PAi}_{\mathrm{op}}
   \le C(R,\Omega) \cdot \left(\frac{\log\log c}{\log c}\right)^{1/2}.
  \label{eq:eps_def}
\end{align}
\end{definition}

\begin{remark}[Sources of the two scales]
$\delta_{\mathrm{Airy}}$ is determined entirely by the Airy kernel $\KAi$:
the Fredholm determinant identity $\det(I-\KAi) = F_{\mathrm{TW}}(0)$
(Tracy--Widom~\cite{TW1994}) gives $\delta_{\mathrm{Airy}} \ge F_{\mathrm{TW}}(0)$.
This is a universal constant of random matrix theory at the GUE edge;
it does not depend on the arithmetic of the sampling sequence.

$\varepsilon_c$ is determined by the \emph{discrepancy}
of the prime sampling sequence $\{p_j^{\mathrm{Ai}}\}$ in the Airy window.
The Montgomery--Vaughan mean square theorem~\cite{MV1974} controls
this discrepancy unconditionally, giving the rate $(\log\log c/\log c)^{1/2}$.
\end{remark}

%%=================================================================
\section{The Perturbative Universality Theorem}
\label{sec:theorem}
%%=================================================================

\begin{theorem}[Perturbative Universality]\label{thm:main}
There exists an explicit $c_0 = c_0(R,\Omega) < \infty$ such that
for all $c \ge c_0$:
\begin{equation}\label{eq:main}
  \inf\sigma(\DAi)
  \ge \delta_{\mathrm{Airy}} - \varepsilon_c
  \ge F_{\mathrm{TW}}(0) - C(R,\Omega)\left(\frac{\log\log c}{\log c}\right)^{1/2}
  > 0.
\end{equation}

In particular:
\begin{enumerate}[\rm (i)]
  \item The spectral gap of $\DAi$ is positive for all large $c$.
  \item The leading term $\delta_{\mathrm{Airy}} \ge 0.96$ is universal:
    it depends only on the Airy kernel, not on the prime distribution.
  \item The arithmetic correction $\varepsilon_c \to 0$:
    the prime structure survives the Airy limit only perturbatively.
  \item The proof is unconditional: no GRH, no conjectures on primes in
    short intervals.
\end{enumerate}
\end{theorem}

\begin{proof}
By the min-max principle:
\begin{align*}
  \ip{\DAi f}{f}
  &= \ip{\PAi(I-\KAi)\PAi f}{f} + \ip{\PAi \Ec \PAi f}{f}\\
  &\ge \delta_{\mathrm{Airy}} \norm{f}^2 - \varepsilon_c \norm{f}^2.
\end{align*}
Insert the bounds \eqref{eq:delta_def} and \eqref{eq:eps_def}.
For $c \ge c_0$ chosen so that $C(R,\Omega)(\log\log c/\log c)^{1/2} < \delta_{\mathrm{Airy}}$,
the right side is $\ge (\delta_{\mathrm{Airy}} - \varepsilon_c)\norm{f}^2 > 0$. \qed
\end{proof}

\begin{remark}[The asymmetry]
The key is the scale separation:
\[
  \frac{\delta_{\mathrm{Airy}}}{\varepsilon_c}
  \sim \frac{0.96}{(\log\log c/\log c)^{1/2}}
  \to +\infty \quad\text{as } c \to \infty.
\]
The universal gap grows (relative to the perturbation) without bound.
The arithmetic perturbation is not merely small --- it becomes \emph{negligible}.
Numerically: the ratio $\mu_A/\mu_K \in [1.88, 4.21]$ for $c \in [100, 10^4]$
(see \texttt{O5\_B\_gap\_persistence.md}).
\end{remark}

%%=================================================================
\section{What the Theorem Says About Universality}
\label{sec:universality}
%%=================================================================

\subsection{The precise universality statement}

The theorem distinguishes sharply between three possible universality regimes:

\begin{center}
\begin{tabular}{lll}
\toprule
Regime & Meaning & Verdict for XXII\\
\midrule
Strong universality & $\inf\sigma(\DAi)$ is \emph{independent} of prime distribution & FALSE\\
Weak universality & $\inf\sigma(\DAi) > 0$ for \emph{any} positive-density sequence & TRUE\\
Perturbative universality & $\inf\sigma(\DAi) = \delta_{\mathrm{Airy}} + O(\varepsilon_c)$,
  $\varepsilon_c \to 0$ & TRUE (Thm.~\ref{thm:main})\\
\bottomrule
\end{tabular}
\end{center}

\medskip

Perturbative universality is the precise statement:
\begin{itemize}
  \item The gap \emph{exists} universally (for any positive-density sequence).
  \item Its \emph{leading value} is universal ($\delta_{\mathrm{Airy}} = $ Airy geometry).
  \item Its \emph{corrections} are arithmetic (prime distribution enters at order $\varepsilon_c$).
  \item The corrections \emph{vanish} as $c \to \infty$.
\end{itemize}

\subsection{Comparison with classical universality}

In random matrix theory, ``edge universality'' means:
the local eigenvalue statistics near the spectral edge of a Wigner matrix
converge to the GUE Tracy--Widom distribution, independent of the
entry distribution.
This is strong universality: the limit is completely distribution-free.

The XXII programme establishes an \emph{operator-level} analogue:
the \emph{spectral gap} of the Airy mismatch operator is
asymptotically distribution-free (prime-distribution-free)
at leading order, with the arithmetic distribution entering
only at vanishing subleading order.

This is perturbative universality: the universal structure
(Airy geometry) is not destroyed by the arithmetic structure
(prime sampling), but dominates it asymptotically.

\subsection{The role of Tracy--Widom}

The constant $F_{\mathrm{TW}}(0) \approx 0.9597$ plays a double role:
\begin{enumerate}
  \item It provides the lower bound on $\delta_{\mathrm{Airy}}$
    (the universal gap is at least $F_{\mathrm{TW}}(0)$).
  \item It provides the threshold $C_{\mathrm{crit}} = F_{\mathrm{TW}}(0)/(\zeta(3/2)-1) \approx 0.595$
    for the Jaffard route (O4/O5 coupling via off-diagonal constant $C$).
\end{enumerate}
In both roles, Tracy--Widom appears as the \emph{quantitative input from universality}
that drives the arithmetic side of the problem.
This is a structural connection: the GUE edge universality constant
directly controls the stability of arithmetic sampling.

%%=================================================================
\section{What the Theorem Does NOT Claim}
\label{sec:notclaim}
%%=================================================================

\begin{remark}[What is NOT claimed]
The theorem does not claim:
\begin{enumerate}
  \item That $\inf\sigma(\DAi)$ is independent of the prime distribution.
    The corrections $\varepsilon_c$ are arithmetic-sensitive
    (sub4 of \texttt{O5\_airy\_gap\_problem.tex}).
  \item That the gap is large for finite $c$.
    For small $c$, $\varepsilon_c$ may exceed $\delta_{\mathrm{Airy}}$;
    the theorem is asymptotic.
  \item Anything about the RH-collapse (O6): whether the gap collapses under RH
    is a separate, finer question that requires O1/O2 (RHP arm).
  \item That O1/O2 (the RHP arm) is unnecessary.
    O1/O2 would upgrade the gap from $c^{-1/2}$ to $c^{-1/3}$ scale
    and give information about the \emph{rate} at which $\varepsilon_c$ decays.
    These are refinements of the main theorem, not prerequisites.
  \item That the proof is optimal. The rate $\varepsilon_c \sim (\log\log c/\log c)^{1/2}$
    from Montgomery--Vaughan is likely not sharp;
    under GRH one would get $\varepsilon_c \sim (\log c)^{-1/2}$.
\end{enumerate}
\end{remark}

%%=================================================================
\section{Proof Dependencies (DAG)}
\label{sec:dag}
%%=================================================================

The proof of Theorem~\ref{thm:main} uses exactly two black-box inputs
from the literature, plus two local lemmas:

\begin{center}
\begin{tabular}{lll}
\toprule
Input & Source & Role\\
\midrule
$F_{\mathrm{TW}}(0) \approx 0.9597$ & Tracy--Widom \cite{TW1994}, Bornemann \cite{Bornemann2010} & $\delta_{\mathrm{Airy}} \ge F_{\mathrm{TW}}$\\
MV mean square & Montgomery--Vaughan \cite{MV1974} & $\varepsilon_c \to 0$\\
\midrule
Bernstein for $\PAi L^2$ & standard FA & Airy-bandlimited $L^\infty$ control\\
PNT counting & Chebyshev / Rosser & $M(c) \sim c^{1/3}/\log c$\\
\bottomrule
\end{tabular}
\end{center}

\medskip
No IIKS structure, no RHP, no Deift-Zhou, no GRH.
The integrable systems machinery (Papers 12--21) provides the
\emph{context} (Airy limit, PSWF convergence) but is not
used in the proof of Theorem~\ref{thm:main} itself.

%%=================================================================
\section{Connection to the Discrete Arm (O4)}
\label{sec:o4}
%%=================================================================

Theorem~\ref{thm:main} (O5 arm) and the Jaffard-Gram bound (O4 arm)
are logically independent. They close P12 by different routes:

\begin{center}
\begin{tabular}{lll}
\toprule
Route & Mechanism & Closes\\
\midrule
O5 (gap-persistence) & $\delta_{\mathrm{Airy}} - \varepsilon_c > 0$ & P12 via Theorem~\ref{thm:main}\\
O4 (Jaffard) & $\lambda_{\min}(G) \ge 1 - C_J$, $C < C_{\mathrm{crit}}$ & P12 via $\lambda_{\min}(A^{\mathrm{arith}}) > \|\KAi\|$\\
\bottomrule
\end{tabular}
\end{center}

\medskip
The O4 route provides quantitative control on $\lambda_{\min}(A^{\mathrm{arith}})$
via the off-diagonal constant $C$.
The O5 route bypasses this and uses $\delta_{\mathrm{Airy}}$ directly.
Both require $C < C_{\mathrm{crit}} = F_{\mathrm{TW}}(0)/(\zeta(3/2)-1) \approx 0.595$
only in the Jaffard route; in the gap-persistence route, $C_{\mathrm{crit}}$ does not appear.

The two routes share a common constant $C$ (the off-diagonal Gram decay):
O4 computes $C$; O5 does not need it.
That O4 was developed first and provides $C$ is a historical artifact
of the programme, not a logical dependency.

\begin{thebibliography}{99}
\bibitem{TW1994}
C.~Tracy and H.~Widom,
\textit{Level-spacing distributions and the Airy kernel},
Comm.\ Math.\ Phys.\ \textbf{159} (1994), 151--174.
\bibitem{Bornemann2010}
F.~Bornemann,
\textit{On the numerical evaluation of Fredholm determinants},
Math.\ Comp.\ \textbf{79} (2010), 871--915.
\bibitem{MV1974}
H.~L.~Montgomery and R.~C.~Vaughan,
\textit{Hilbert's inequality},
J.\ London Math.\ Soc.\ (2) \textbf{8} (1974), 73--82.
\bibitem{MV_book}
H.~L.~Montgomery and R.~C.~Vaughan,
\textit{Multiplicative Number Theory I},
Cambridge University Press, 2006.
\bibitem{Beurling1938}
A.~Beurling,
\textit{Sur les int\'egrales de Fourier absolument convergentes},
IX\`eme Congr\`es Math.\ Scand., 1938.
\bibitem{PaperXXII}
S.~Schmalnauer,
\textit{Edge Coercivity and the Airy Boundary Operator} (Paper~XXII),
preprint, 2026.
\end{thebibliography}

\end{document}
